% ==============================================================================
% PLANTILLA_ALUMNO.TEX — MACRO Y DEFINICIÓN DE FICHA PARA ALUMNADO (NEAE)
% Asignatura: Metodologías de Desarrollo Ágil (MDA) — Universidad de Granada
% Proyecto AUTOVIDA (26-175) | Reto CeRA — Agenda Accesible PTVAL
% ==============================================================================
% Campos canónicos: 9 (Sin Occupation, Education ni Company Information)
% Uso: % ==============================================================================
% PLANTILLA_ALUMNO.TEX — MACRO Y DEFINICIÓN DE FICHA PARA ALUMNADO (NEAE)
% Asignatura: Metodologías de Desarrollo Ágil (MDA) — Universidad de Granada
% Proyecto AUTOVIDA (26-175) | Reto CeRA — Agenda Accesible PTVAL
% ==============================================================================
% Campos canónicos: 9 (Sin Occupation, Education ni Company Information)
% Uso: % ==============================================================================
% PLANTILLA_ALUMNO.TEX — MACRO Y DEFINICIÓN DE FICHA PARA ALUMNADO (NEAE)
% Asignatura: Metodologías de Desarrollo Ágil (MDA) — Universidad de Granada
% Proyecto AUTOVIDA (26-175) | Reto CeRA — Agenda Accesible PTVAL
% ==============================================================================
% Campos canónicos: 9 (Sin Occupation, Education ni Company Information)
% Uso: % ==============================================================================
% PLANTILLA_ALUMNO.TEX — MACRO Y DEFINICIÓN DE FICHA PARA ALUMNADO (NEAE)
% Asignatura: Metodologías de Desarrollo Ágil (MDA) — Universidad de Granada
% Proyecto AUTOVIDA (26-175) | Reto CeRA — Agenda Accesible PTVAL
% ==============================================================================
% Campos canónicos: 9 (Sin Occupation, Education ni Company Information)
% Uso: \input{plantilla_alumno.tex} y luego invocar \fichaAlumno[...] o \fichaAlumnoBlanco
% ==============================================================================

\ifdefined\isSubFile
    % Si se incluye desde otro documento, los paquetes ya están cargados
\else
    \documentclass[10pt]{article}
    \usepackage[utf8]{inputenc}
    \usepackage[spanish,es-tabla]{babel}
    \usepackage[margin=0.8cm,top=0.6cm,bottom=0.6cm]{geometry}
    \usepackage{helvet}
    \renewcommand{\familydefault}{\sfdefault}
    \usepackage{microtype}
    \usepackage{graphicx}
    \usepackage{xcolor}
    \usepackage{tikz}
    \usepackage{enumitem}
    \usepackage{fontawesome5}
    \usepackage{tcolorbox}
    \tcbuselibrary{skins}
    \pagestyle{empty}
\fi

% Colores del sistema
\providecolor{primaryColor}{RGB}{20, 65, 105}     % Azul corporativo
\providecolor{secondaryColor}{RGB}{35, 115, 170}  % Azul medio accesible
\providecolor{tealColor}{RGB}{20, 130, 140}       % Verde azulado
\providecolor{greenColor}{RGB}{30, 135, 75}       % Verde bosque
\providecolor{accentColor}{RGB}{215, 105, 25}     % Naranja cálido
\providecolor{purpleColor}{RGB}{120, 60, 150}     % Púrpura suave
\providecolor{dangerColor}{RGB}{185, 45, 35}      % Rojo criticidad
\providecolor{bgCard}{RGB}{252, 253, 255}         % Fondo limpio
\providecolor{borderCard}{RGB}{215, 225, 235}     % Borde sutil
\providecolor{textDark}{RGB}{30, 40, 50}          % Texto principal oscuro
\providecolor{textMuted}{RGB}{95, 110, 125}       % Texto secundario

% Configuración de listas
\setlist[itemize]{leftmargin=3.0mm, itemsep=0.4pt, parsep=0pt, topsep=0.4pt, partopsep=0pt}

% Estilos de caja fija para Alumno
\tcbset{
    cardAlumnoBox/.style={
        enhanced,
        boxrule=0.65pt,
        arc=1.8mm,
        height=32.0mm,
        left=5.0pt, right=5.0pt, top=1.8pt, bottom=1.8pt,
        boxsep=0pt,
        colback=bgCard,
        colframe=borderCard,
        fonttitle=\bfseries\small,
        fontupper=\fontsize{8.5pt}{10.0pt}\selectfont,
    },
    cardAlumnoWideBox/.style={
        enhanced,
        boxrule=0.65pt,
        arc=1.8mm,
        height=22.5mm,
        left=6.0pt, right=6.0pt, top=2.0pt, bottom=2.0pt,
        boxsep=0pt,
        colback=bgCard,
        colframe=borderCard,
        fonttitle=\bfseries\small,
    }
}

% Macro para badges visuales (si no existe)
\providecommand{\badge}[2]{%
    \tikz[baseline=(b.base)] \node[fill=#1!12, text=#1!90!black, draw=#1!35, line width=0.4pt, rounded corners=1.2mm, inner sep=2.0pt, font=\bfseries\scriptsize] (b) {#2};%
}

% Definición de claves con pgfkeys para \fichaAlumno
\pgfkeys{
    /fichaAlumno/.is family, /fichaAlumno,
    default/.style = {
        bannerTitulo = {\faUserGraduate\ FICHA DE PERSONA: ALUMNADO (METODOLOGÍA DCU)},
        subtitulo = {Proyecto AUTOVIDA (26-175) \textbar\ Reto CeRA — Agenda Accesible PTVAL \textbar\ Asignatura MDA (UGR)},
        equipo = {\textbf{Equipo CeRA-MDA} (5 integrantes)\\\ Course 2026/2027 \textbar\ \textbf{Práctica 1}},
        nombre = {[Nombre y Apellidos del Alumno]},
        nombreCorto = {[ALUMNO]},
        edad = {[Edad]},
        icono = {\faUserGraduate},
        foto = {},
        badgeRol = {[Perfil NEAE / Aula]},
        badgeTarea = {Tareas Rotativas},
        centro = {CEE Purísima Concepción},
        cita = {[Cita textual representativa que exprese cómo vive su rutina, qué le ayuda a anticipar y qué le genera ansiedad o rechazo].},
        demographic = {
            \item \textbf{Edad y Género (Age / Gender):} [[Edad] años, [Género]].
            \item \textbf{Ubicación (Location):} [[Granada capital / Municipio]].
            \item \textbf{Perfil diagnóstico:} [[Diagnóstico cognitivo, motor o sensorial]].
        },
        computerExp = {
            \item \textbf{Habilidades informáticas:} [[Nivel táctil: autónomo, asistido]].
            \item \textbf{Años de experiencia:} [[Años de contacto con tablets]].
            \item \textbf{Hábitos de uso:} [[Apps frecuentes: vídeos, pictogramas, juegos]].
            \item \textbf{Autonomía digital:} [[Baja / Media / Asistida]].
        },
        productExp = {
            \item \textbf{Productos similares:} [[PECS en papel, velcro, apps previas]].
            \item \textbf{Limitaciones actuales:} [[Pérdida de material, rigidez, frustración]].
            \item \textbf{Tendencias de uso:} [[Receptividad a imágenes reales vs pictogramas]].
        },
        tasks = {
            \item \textbf{Tareas Primarias:} 1) [[Login accesible]] $\rightarrow$ 2) [[Ver tarea actual]] $\rightarrow$ 3) [[Ejecutar pasos guiados]] $\rightarrow$ 4) [[Confirmar finalización]].
            \item \textbf{Tareas Secundarias:} [[Ver siguiente tarea / pedir ayuda]].
        },
        domainKnowledge = {
            \item \textbf{Comprensión del área:} [[Orientación en el colegio, aulas y profesores]].
            \item \textbf{Comprensión temporal:} [[Visual (disco/reloj de arena), no numérica]].
            \item \textbf{Destreza procedimental:} [[Grado de autonomía en la tarea asignada]].
        },
        technology = {
            \item \textbf{Hardware:} [[Tablet Android / pulsador externo]].
            \item \textbf{Software:} [[Android / síntesis de voz]].
            \item \textbf{Herramientas habituales:} [[Conexión Wi-Fi / atril]].
        },
        attitudes = {
            \item \textbf{Preferencias:} [[Refuerzos sensoriales favoritos, limpieza visual]].
            \item \textbf{Miedo a la tecnología:} [[Reacción a errores, sonidos o bloqueos]].
            \item \textbf{Valores:} [[Sentimiento de logro, orgullo por autonomía]].
        },
        learningStyle = {
            \item \textbf{Perfil de aprendizaje:} [[\textbf{Visual} / \textbf{Auditivo} / \textbf{Kinestésico}]].
            \item \textbf{Canales óptimos:} [[Fotos reales / Pictogramas ARASAAC / Vídeos]].
            \item \textbf{Metodología:} [[Modelado guiado, repetición, ensayo y éxito]].
        },
        criticality = {
            \item \textbf{Consecuencias de un error:} \textit{Emocional:} [[¿Bloqueo, crisis de frustración, rechazo?]] \textit{Operativo:} [[¿Desorientación, interrupción de la tarea?]]
            \item \textbf{Tolerancia del software:} [[Acciones no destructivas, confirmación siempre asistida, botón SOS visible]].
        },
        escenario = {[[Narrar en 3-4 líneas una situación real: a qué hora usa la agenda, cómo anticipa su rutina con apoyos visuales y cómo supera el momento crítico logrando autonomía sin estrés.]]},
        badgesAjustes = {
            \badge{primaryColor}{\faKey\ Login: [[Método]]} \hspace{2mm}
            \badge{secondaryColor}{\faIcons\ [[ARASAAC / Fotos]]} \hspace{2mm}
            \badge{accentColor}{\faClock\ [[Temporizador]]} \hspace{2mm}
            \badge{purpleColor}{\faVolumeUp\ [[Voz TTS]]}%
        }
    },
    bannerTitulo/.store in = \fAlBannerTitulo,
    subtitulo/.store in = \fAlSubtitulo,
    equipo/.store in = \fAlEquipo,
    nombre/.store in = \fAlNombre,
    nombreCorto/.store in = \fAlNombreCorto,
    edad/.store in = \fAlEdad,
    icono/.store in = \fAlIcono,
    foto/.store in = \fAlFoto,
    badgeRol/.store in = \fAlBadgeRol,
    badgeTarea/.store in = \fAlBadgeTarea,
    centro/.store in = \fAlCentro,
    cita/.store in = \fAlCita,
    demographic/.store in = \fAlDemographic,
    computerExp/.store in = \fAlComputerExp,
    productExp/.store in = \fAlProductExp,
    tasks/.store in = \fAlTasks,
    domainKnowledge/.store in = \fAlDomainKnowledge,
    technology/.store in = \fAlTechnology,
    attitudes/.store in = \fAlAttitudes,
    learningStyle/.store in = \fAlLearningStyle,
    criticality/.store in = \fAlCriticality,
    escenario/.store in = \fAlEscenario,
    badgesAjustes/.store in = \fAlBadgesAjustes,
}

% COMANDO PRINCIPAL \fichaAlumno
\newcommand{\fichaAlumno}[1][]{%
    \begingroup
    \pgfkeys{/fichaAlumno, default, #1}%
    % 1. BANNER SUPERIOR
    \begin{tcolorbox}[
        enhanced,
        colback=primaryColor,
        colframe=primaryColor,
        arc=1.6mm,
        left=7pt, right=7pt, top=2.5pt, bottom=2.5pt,
        boxsep=0pt, coltext=white
    ]
        \begin{minipage}[c]{0.70\textwidth}
            {\normalsize\textbf{\fAlBannerTitulo}}\\[0.2mm]
            {\tiny \fAlSubtitulo}
        \end{minipage}\hfill
        \begin{minipage}[c]{0.29\textwidth}
            \raggedleft\tiny
            \fAlEquipo
        \end{minipage}
    \end{tcolorbox}
    \vspace{0.8mm}
    % 2. TARJETA DE IDENTIDAD DEL ALUMNO
    \begin{tcolorbox}[
        enhanced,
        colback=white,
        colframe=primaryColor!65,
        arc=2.0mm,
        boxrule=0.8pt,
        left=6pt, right=6pt, top=3.0pt, bottom=3.0pt,
        boxsep=0pt,
        drop shadow=black!5
    ]
        \begin{minipage}[c]{0.10\textwidth}
            \centering
            \begin{tikzpicture}
                \ifx\fAlFoto\empty
                    \draw[fill=primaryColor!8, draw=primaryColor!30, line width=0.5pt, rounded corners=2mm] (0,0) rectangle (1.4, 1.55);
                    \node at (0.7, 0.95) {\color{primaryColor}\scalebox{1.5}{\fAlIcono}};
                    \node[align=center] at (0.7, 0.28) {\tiny\textbf{\color{primaryColor}\fAlNombreCorto}\\\fontsize{4pt}{5pt}\selectfont\color{textMuted}\fAlEdad};
                \else
                    \begin{scope}
                        \clip[rounded corners=2mm] (0,0) rectangle (1.45, 1.55);
                        \node[inner sep=0pt] at (0.725, 0.775) {\includegraphics[width=1.45cm,height=1.55cm,keepaspectratio=false]{\fAlFoto}};
                    \end{scope}
                    \draw[draw=primaryColor!50, line width=0.8pt, rounded corners=2mm] (0,0) rectangle (1.45, 1.55);
                \fi
            \end{tikzpicture}
        \end{minipage}\hfill
        \begin{minipage}[c]{0.88\textwidth}
            \begin{minipage}[b]{0.68\textwidth}\raggedright
                {\large\textbf{\color{primaryColor}\fAlNombre}} \hspace{2mm}
                \badge{secondaryColor}{\fAlBadgeRol} \badge{accentColor}{\fAlBadgeTarea}
            \end{minipage}\hfill
            \begin{minipage}[b]{0.31\textwidth}
                \raggedleft\small\textbf{\color{primaryColor}\faSchool\ \fAlCentro}
            \end{minipage}\\[0.6mm]
            \tikz[baseline] \node[fill=primaryColor!5, draw=primaryColor!25, rounded corners=1.2mm, inner sep=2.5pt, text width=0.98\linewidth] {
                \small\itshape\color{primaryColor}
                \faQuoteLeft\ ``\fAlCita''\ \faQuoteRight
            };
        \end{minipage}
    \end{tcolorbox}
    \vspace{0.8mm}
    % 3. LAS 9 CAJAS DE ALUMNOS (4 FILAS X 2 COLS + 1 FILA COMPLETA)
    \noindent\begin{minipage}[t]{0.485\textwidth}
        \begin{tcolorbox}[cardAlumnoBox, colframe=primaryColor!50, coltitle=primaryColor, title={\faIdCard\ 1. Demographic characteristics}]
            \fontsize{8.2pt}{9.8pt}\selectfont
            \begin{itemize}
                \fAlDemographic
            \end{itemize}
        \end{tcolorbox}
    \end{minipage}\hfill
    \begin{minipage}[t]{0.485\textwidth}
        \begin{tcolorbox}[cardAlumnoBox, colframe=secondaryColor!60, coltitle=secondaryColor!90!black, title={\faLaptop\ 2. Computer experience}]
            \fontsize{8.2pt}{9.8pt}\selectfont
            \begin{itemize}
                \fAlComputerExp
            \end{itemize}
        \end{tcolorbox}
    \end{minipage}
    \vspace{0.8mm}
    \noindent\begin{minipage}[t]{0.485\textwidth}
        \begin{tcolorbox}[cardAlumnoBox, colframe=secondaryColor!60, coltitle=secondaryColor!90!black, title={\faSearchPlus\ 3. Specific product experience}]
            \fontsize{8.2pt}{9.8pt}\selectfont
            \begin{itemize}
                \fAlProductExp
            \end{itemize}
        \end{tcolorbox}
    \end{minipage}\hfill
    \begin{minipage}[t]{0.485\textwidth}
        \begin{tcolorbox}[cardAlumnoBox, colframe=greenColor!65, coltitle=greenColor!90!black, title={\faTasks\ 4. Tasks (Primary and secondary)}]
            \fontsize{8.2pt}{9.8pt}\selectfont
            \begin{itemize}
                \fAlTasks
            \end{itemize}
        \end{tcolorbox}
    \end{minipage}
    \vspace{0.8mm}
    \noindent\begin{minipage}[t]{0.485\textwidth}
        \begin{tcolorbox}[cardAlumnoBox, colframe=accentColor!65, coltitle=accentColor!90!black, title={\faBrain\ 5. Domain knowledge}]
            \fontsize{8.2pt}{9.8pt}\selectfont
            \begin{itemize}
                \fAlDomainKnowledge
            \end{itemize}
        \end{tcolorbox}
    \end{minipage}\hfill
    \begin{minipage}[t]{0.485\textwidth}
        \begin{tcolorbox}[cardAlumnoBox, colframe=tealColor!60, coltitle=tealColor!90!black, title={\faTools\ 6. Technology available}]
            \fontsize{8.2pt}{9.8pt}\selectfont
            \begin{itemize}
                \fAlTechnology
            \end{itemize}
        \end{tcolorbox}
    \end{minipage}
    \vspace{0.8mm}
    \noindent\begin{minipage}[t]{0.485\textwidth}
        \begin{tcolorbox}[cardAlumnoBox, colframe=purpleColor!60, coltitle=purpleColor!90!black, title={\faHeart\ 7. Attitudes and values}]
            \fontsize{8.2pt}{9.8pt}\selectfont
            \begin{itemize}
                \fAlAttitudes
            \end{itemize}
        \end{tcolorbox}
    \end{minipage}\hfill
    \begin{minipage}[t]{0.485\textwidth}
        \begin{tcolorbox}[cardAlumnoBox, colframe=purpleColor!60, coltitle=purpleColor!90!black, title={\faEye\ 8. Learning style}]
            \fontsize{8.2pt}{9.8pt}\selectfont
            \begin{itemize}
                \fAlLearningStyle
            \end{itemize}
        \end{tcolorbox}
    \end{minipage}
    \vspace{0.8mm}
    % FILA 5: Punto 9 (ANCHO COMPLETO — CRITICALITY OF ERRORS)
    \noindent\begin{minipage}[t]{\textwidth}
        \begin{tcolorbox}[cardAlumnoWideBox, colframe=dangerColor!65, coltitle=dangerColor!90!black, title={\faExclamationTriangle\ 9. Criticality of errors (Consecuencias de un error y tolerancia del software)}]
            \fontsize{8.2pt}{9.8pt}\selectfont
            \begin{itemize}
                \fAlCriticality
            \end{itemize}
        \end{tcolorbox}
    \end{minipage}
    \vspace{0.8mm}
    % 4. BARRA INFERIOR DE INTEGRACIÓN CeRA
    \begin{tcolorbox}[
        enhanced,
        colback=primaryColor!4,
        colframe=primaryColor!45,
        arc=1.6mm,
        left=6pt, right=6pt, top=2.5pt, bottom=2.5pt,
        boxsep=0pt
    ]
        \scriptsize
        \textbf{\color{primaryColor}\faMapSigns\ Escenario de Uso Contextual (Día típico):} \fAlEscenario\\
        \vspace{0.3mm}
        \hrule height 0.35pt \relax
        \vspace{0.5mm}
        \centerline{%
            \textbf{\color{primaryColor}\faSlidersH\ Ajustes UI Accesible (CeRA):} \hspace{2mm}
            \fAlBadgesAjustes%
        }
    \end{tcolorbox}
    \newpage
    \endgroup
}

% COMANDO PARA PLANTILLA EN BLANCO DE ALUMNADO
\newcommand{\fichaAlumnoBlanco}{%
    \fichaAlumno[%
        bannerTitulo={\faUserEdit\ FICHA DE PERSONA: ALUMNADO \textbar\ PLANTILLA DE TRABAJO},
        cita={[Cita textual que exprese cómo vive su rutina, qué le ayuda a anticipar y qué le genera ansiedad o rechazo].}%
    ]%
}

\ifdefined\isSubFile
    \expandafter\endinput
\fi

\begin{document}
\fichaAlumnoBlanco
\end{document}
 y luego invocar \fichaAlumno[...] o \fichaAlumnoBlanco
% ==============================================================================

\ifdefined\isSubFile
    % Si se incluye desde otro documento, los paquetes ya están cargados
\else
    \documentclass[10pt]{article}
    \usepackage[utf8]{inputenc}
    \usepackage[spanish,es-tabla]{babel}
    \usepackage[margin=0.8cm,top=0.6cm,bottom=0.6cm]{geometry}
    \usepackage{helvet}
    \renewcommand{\familydefault}{\sfdefault}
    \usepackage{microtype}
    \usepackage{graphicx}
    \usepackage{xcolor}
    \usepackage{tikz}
    \usepackage{enumitem}
    \usepackage{fontawesome5}
    \usepackage{tcolorbox}
    \tcbuselibrary{skins}
    \pagestyle{empty}
\fi

% Colores del sistema
\providecolor{primaryColor}{RGB}{20, 65, 105}     % Azul corporativo
\providecolor{secondaryColor}{RGB}{35, 115, 170}  % Azul medio accesible
\providecolor{tealColor}{RGB}{20, 130, 140}       % Verde azulado
\providecolor{greenColor}{RGB}{30, 135, 75}       % Verde bosque
\providecolor{accentColor}{RGB}{215, 105, 25}     % Naranja cálido
\providecolor{purpleColor}{RGB}{120, 60, 150}     % Púrpura suave
\providecolor{dangerColor}{RGB}{185, 45, 35}      % Rojo criticidad
\providecolor{bgCard}{RGB}{252, 253, 255}         % Fondo limpio
\providecolor{borderCard}{RGB}{215, 225, 235}     % Borde sutil
\providecolor{textDark}{RGB}{30, 40, 50}          % Texto principal oscuro
\providecolor{textMuted}{RGB}{95, 110, 125}       % Texto secundario

% Configuración de listas
\setlist[itemize]{leftmargin=3.0mm, itemsep=0.4pt, parsep=0pt, topsep=0.4pt, partopsep=0pt}

% Estilos de caja fija para Alumno
\tcbset{
    cardAlumnoBox/.style={
        enhanced,
        boxrule=0.65pt,
        arc=1.8mm,
        height=32.0mm,
        left=5.0pt, right=5.0pt, top=1.8pt, bottom=1.8pt,
        boxsep=0pt,
        colback=bgCard,
        colframe=borderCard,
        fonttitle=\bfseries\small,
        fontupper=\fontsize{8.5pt}{10.0pt}\selectfont,
    },
    cardAlumnoWideBox/.style={
        enhanced,
        boxrule=0.65pt,
        arc=1.8mm,
        height=22.5mm,
        left=6.0pt, right=6.0pt, top=2.0pt, bottom=2.0pt,
        boxsep=0pt,
        colback=bgCard,
        colframe=borderCard,
        fonttitle=\bfseries\small,
    }
}

% Macro para badges visuales (si no existe)
\providecommand{\badge}[2]{%
    \tikz[baseline=(b.base)] \node[fill=#1!12, text=#1!90!black, draw=#1!35, line width=0.4pt, rounded corners=1.2mm, inner sep=2.0pt, font=\bfseries\scriptsize] (b) {#2};%
}

% Definición de claves con pgfkeys para \fichaAlumno
\pgfkeys{
    /fichaAlumno/.is family, /fichaAlumno,
    default/.style = {
        bannerTitulo = {\faUserGraduate\ FICHA DE PERSONA: ALUMNADO (METODOLOGÍA DCU)},
        subtitulo = {Proyecto AUTOVIDA (26-175) \textbar\ Reto CeRA — Agenda Accesible PTVAL \textbar\ Asignatura MDA (UGR)},
        equipo = {\textbf{Equipo CeRA-MDA} (5 integrantes)\\\ Course 2026/2027 \textbar\ \textbf{Práctica 1}},
        nombre = {[Nombre y Apellidos del Alumno]},
        nombreCorto = {[ALUMNO]},
        edad = {[Edad]},
        icono = {\faUserGraduate},
        foto = {},
        badgeRol = {[Perfil NEAE / Aula]},
        badgeTarea = {Tareas Rotativas},
        centro = {CEE Purísima Concepción},
        cita = {[Cita textual representativa que exprese cómo vive su rutina, qué le ayuda a anticipar y qué le genera ansiedad o rechazo].},
        demographic = {
            \item \textbf{Edad y Género (Age / Gender):} [[Edad] años, [Género]].
            \item \textbf{Ubicación (Location):} [[Granada capital / Municipio]].
            \item \textbf{Perfil diagnóstico:} [[Diagnóstico cognitivo, motor o sensorial]].
        },
        computerExp = {
            \item \textbf{Habilidades informáticas:} [[Nivel táctil: autónomo, asistido]].
            \item \textbf{Años de experiencia:} [[Años de contacto con tablets]].
            \item \textbf{Hábitos de uso:} [[Apps frecuentes: vídeos, pictogramas, juegos]].
            \item \textbf{Autonomía digital:} [[Baja / Media / Asistida]].
        },
        productExp = {
            \item \textbf{Productos similares:} [[PECS en papel, velcro, apps previas]].
            \item \textbf{Limitaciones actuales:} [[Pérdida de material, rigidez, frustración]].
            \item \textbf{Tendencias de uso:} [[Receptividad a imágenes reales vs pictogramas]].
        },
        tasks = {
            \item \textbf{Tareas Primarias:} 1) [[Login accesible]] $\rightarrow$ 2) [[Ver tarea actual]] $\rightarrow$ 3) [[Ejecutar pasos guiados]] $\rightarrow$ 4) [[Confirmar finalización]].
            \item \textbf{Tareas Secundarias:} [[Ver siguiente tarea / pedir ayuda]].
        },
        domainKnowledge = {
            \item \textbf{Comprensión del área:} [[Orientación en el colegio, aulas y profesores]].
            \item \textbf{Comprensión temporal:} [[Visual (disco/reloj de arena), no numérica]].
            \item \textbf{Destreza procedimental:} [[Grado de autonomía en la tarea asignada]].
        },
        technology = {
            \item \textbf{Hardware:} [[Tablet Android / pulsador externo]].
            \item \textbf{Software:} [[Android / síntesis de voz]].
            \item \textbf{Herramientas habituales:} [[Conexión Wi-Fi / atril]].
        },
        attitudes = {
            \item \textbf{Preferencias:} [[Refuerzos sensoriales favoritos, limpieza visual]].
            \item \textbf{Miedo a la tecnología:} [[Reacción a errores, sonidos o bloqueos]].
            \item \textbf{Valores:} [[Sentimiento de logro, orgullo por autonomía]].
        },
        learningStyle = {
            \item \textbf{Perfil de aprendizaje:} [[\textbf{Visual} / \textbf{Auditivo} / \textbf{Kinestésico}]].
            \item \textbf{Canales óptimos:} [[Fotos reales / Pictogramas ARASAAC / Vídeos]].
            \item \textbf{Metodología:} [[Modelado guiado, repetición, ensayo y éxito]].
        },
        criticality = {
            \item \textbf{Consecuencias de un error:} \textit{Emocional:} [[¿Bloqueo, crisis de frustración, rechazo?]] \textit{Operativo:} [[¿Desorientación, interrupción de la tarea?]]
            \item \textbf{Tolerancia del software:} [[Acciones no destructivas, confirmación siempre asistida, botón SOS visible]].
        },
        escenario = {[[Narrar en 3-4 líneas una situación real: a qué hora usa la agenda, cómo anticipa su rutina con apoyos visuales y cómo supera el momento crítico logrando autonomía sin estrés.]]},
        badgesAjustes = {
            \badge{primaryColor}{\faKey\ Login: [[Método]]} \hspace{2mm}
            \badge{secondaryColor}{\faIcons\ [[ARASAAC / Fotos]]} \hspace{2mm}
            \badge{accentColor}{\faClock\ [[Temporizador]]} \hspace{2mm}
            \badge{purpleColor}{\faVolumeUp\ [[Voz TTS]]}%
        }
    },
    bannerTitulo/.store in = \fAlBannerTitulo,
    subtitulo/.store in = \fAlSubtitulo,
    equipo/.store in = \fAlEquipo,
    nombre/.store in = \fAlNombre,
    nombreCorto/.store in = \fAlNombreCorto,
    edad/.store in = \fAlEdad,
    icono/.store in = \fAlIcono,
    foto/.store in = \fAlFoto,
    badgeRol/.store in = \fAlBadgeRol,
    badgeTarea/.store in = \fAlBadgeTarea,
    centro/.store in = \fAlCentro,
    cita/.store in = \fAlCita,
    demographic/.store in = \fAlDemographic,
    computerExp/.store in = \fAlComputerExp,
    productExp/.store in = \fAlProductExp,
    tasks/.store in = \fAlTasks,
    domainKnowledge/.store in = \fAlDomainKnowledge,
    technology/.store in = \fAlTechnology,
    attitudes/.store in = \fAlAttitudes,
    learningStyle/.store in = \fAlLearningStyle,
    criticality/.store in = \fAlCriticality,
    escenario/.store in = \fAlEscenario,
    badgesAjustes/.store in = \fAlBadgesAjustes,
}

% COMANDO PRINCIPAL \fichaAlumno
\newcommand{\fichaAlumno}[1][]{%
    \begingroup
    \pgfkeys{/fichaAlumno, default, #1}%
    % 1. BANNER SUPERIOR
    \begin{tcolorbox}[
        enhanced,
        colback=primaryColor,
        colframe=primaryColor,
        arc=1.6mm,
        left=7pt, right=7pt, top=2.5pt, bottom=2.5pt,
        boxsep=0pt, coltext=white
    ]
        \begin{minipage}[c]{0.70\textwidth}
            {\normalsize\textbf{\fAlBannerTitulo}}\\[0.2mm]
            {\tiny \fAlSubtitulo}
        \end{minipage}\hfill
        \begin{minipage}[c]{0.29\textwidth}
            \raggedleft\tiny
            \fAlEquipo
        \end{minipage}
    \end{tcolorbox}
    \vspace{0.8mm}
    % 2. TARJETA DE IDENTIDAD DEL ALUMNO
    \begin{tcolorbox}[
        enhanced,
        colback=white,
        colframe=primaryColor!65,
        arc=2.0mm,
        boxrule=0.8pt,
        left=6pt, right=6pt, top=3.0pt, bottom=3.0pt,
        boxsep=0pt,
        drop shadow=black!5
    ]
        \begin{minipage}[c]{0.10\textwidth}
            \centering
            \begin{tikzpicture}
                \ifx\fAlFoto\empty
                    \draw[fill=primaryColor!8, draw=primaryColor!30, line width=0.5pt, rounded corners=2mm] (0,0) rectangle (1.4, 1.55);
                    \node at (0.7, 0.95) {\color{primaryColor}\scalebox{1.5}{\fAlIcono}};
                    \node[align=center] at (0.7, 0.28) {\tiny\textbf{\color{primaryColor}\fAlNombreCorto}\\\fontsize{4pt}{5pt}\selectfont\color{textMuted}\fAlEdad};
                \else
                    \begin{scope}
                        \clip[rounded corners=2mm] (0,0) rectangle (1.45, 1.55);
                        \node[inner sep=0pt] at (0.725, 0.775) {\includegraphics[width=1.45cm,height=1.55cm,keepaspectratio=false]{\fAlFoto}};
                    \end{scope}
                    \draw[draw=primaryColor!50, line width=0.8pt, rounded corners=2mm] (0,0) rectangle (1.45, 1.55);
                \fi
            \end{tikzpicture}
        \end{minipage}\hfill
        \begin{minipage}[c]{0.88\textwidth}
            \begin{minipage}[b]{0.68\textwidth}\raggedright
                {\large\textbf{\color{primaryColor}\fAlNombre}} \hspace{2mm}
                \badge{secondaryColor}{\fAlBadgeRol} \badge{accentColor}{\fAlBadgeTarea}
            \end{minipage}\hfill
            \begin{minipage}[b]{0.31\textwidth}
                \raggedleft\small\textbf{\color{primaryColor}\faSchool\ \fAlCentro}
            \end{minipage}\\[0.6mm]
            \tikz[baseline] \node[fill=primaryColor!5, draw=primaryColor!25, rounded corners=1.2mm, inner sep=2.5pt, text width=0.98\linewidth] {
                \small\itshape\color{primaryColor}
                \faQuoteLeft\ ``\fAlCita''\ \faQuoteRight
            };
        \end{minipage}
    \end{tcolorbox}
    \vspace{0.8mm}
    % 3. LAS 9 CAJAS DE ALUMNOS (4 FILAS X 2 COLS + 1 FILA COMPLETA)
    \noindent\begin{minipage}[t]{0.485\textwidth}
        \begin{tcolorbox}[cardAlumnoBox, colframe=primaryColor!50, coltitle=primaryColor, title={\faIdCard\ 1. Demographic characteristics}]
            \fontsize{8.2pt}{9.8pt}\selectfont
            \begin{itemize}
                \fAlDemographic
            \end{itemize}
        \end{tcolorbox}
    \end{minipage}\hfill
    \begin{minipage}[t]{0.485\textwidth}
        \begin{tcolorbox}[cardAlumnoBox, colframe=secondaryColor!60, coltitle=secondaryColor!90!black, title={\faLaptop\ 2. Computer experience}]
            \fontsize{8.2pt}{9.8pt}\selectfont
            \begin{itemize}
                \fAlComputerExp
            \end{itemize}
        \end{tcolorbox}
    \end{minipage}
    \vspace{0.8mm}
    \noindent\begin{minipage}[t]{0.485\textwidth}
        \begin{tcolorbox}[cardAlumnoBox, colframe=secondaryColor!60, coltitle=secondaryColor!90!black, title={\faSearchPlus\ 3. Specific product experience}]
            \fontsize{8.2pt}{9.8pt}\selectfont
            \begin{itemize}
                \fAlProductExp
            \end{itemize}
        \end{tcolorbox}
    \end{minipage}\hfill
    \begin{minipage}[t]{0.485\textwidth}
        \begin{tcolorbox}[cardAlumnoBox, colframe=greenColor!65, coltitle=greenColor!90!black, title={\faTasks\ 4. Tasks (Primary and secondary)}]
            \fontsize{8.2pt}{9.8pt}\selectfont
            \begin{itemize}
                \fAlTasks
            \end{itemize}
        \end{tcolorbox}
    \end{minipage}
    \vspace{0.8mm}
    \noindent\begin{minipage}[t]{0.485\textwidth}
        \begin{tcolorbox}[cardAlumnoBox, colframe=accentColor!65, coltitle=accentColor!90!black, title={\faBrain\ 5. Domain knowledge}]
            \fontsize{8.2pt}{9.8pt}\selectfont
            \begin{itemize}
                \fAlDomainKnowledge
            \end{itemize}
        \end{tcolorbox}
    \end{minipage}\hfill
    \begin{minipage}[t]{0.485\textwidth}
        \begin{tcolorbox}[cardAlumnoBox, colframe=tealColor!60, coltitle=tealColor!90!black, title={\faTools\ 6. Technology available}]
            \fontsize{8.2pt}{9.8pt}\selectfont
            \begin{itemize}
                \fAlTechnology
            \end{itemize}
        \end{tcolorbox}
    \end{minipage}
    \vspace{0.8mm}
    \noindent\begin{minipage}[t]{0.485\textwidth}
        \begin{tcolorbox}[cardAlumnoBox, colframe=purpleColor!60, coltitle=purpleColor!90!black, title={\faHeart\ 7. Attitudes and values}]
            \fontsize{8.2pt}{9.8pt}\selectfont
            \begin{itemize}
                \fAlAttitudes
            \end{itemize}
        \end{tcolorbox}
    \end{minipage}\hfill
    \begin{minipage}[t]{0.485\textwidth}
        \begin{tcolorbox}[cardAlumnoBox, colframe=purpleColor!60, coltitle=purpleColor!90!black, title={\faEye\ 8. Learning style}]
            \fontsize{8.2pt}{9.8pt}\selectfont
            \begin{itemize}
                \fAlLearningStyle
            \end{itemize}
        \end{tcolorbox}
    \end{minipage}
    \vspace{0.8mm}
    % FILA 5: Punto 9 (ANCHO COMPLETO — CRITICALITY OF ERRORS)
    \noindent\begin{minipage}[t]{\textwidth}
        \begin{tcolorbox}[cardAlumnoWideBox, colframe=dangerColor!65, coltitle=dangerColor!90!black, title={\faExclamationTriangle\ 9. Criticality of errors (Consecuencias de un error y tolerancia del software)}]
            \fontsize{8.2pt}{9.8pt}\selectfont
            \begin{itemize}
                \fAlCriticality
            \end{itemize}
        \end{tcolorbox}
    \end{minipage}
    \vspace{0.8mm}
    % 4. BARRA INFERIOR DE INTEGRACIÓN CeRA
    \begin{tcolorbox}[
        enhanced,
        colback=primaryColor!4,
        colframe=primaryColor!45,
        arc=1.6mm,
        left=6pt, right=6pt, top=2.5pt, bottom=2.5pt,
        boxsep=0pt
    ]
        \scriptsize
        \textbf{\color{primaryColor}\faMapSigns\ Escenario de Uso Contextual (Día típico):} \fAlEscenario\\
        \vspace{0.3mm}
        \hrule height 0.35pt \relax
        \vspace{0.5mm}
        \centerline{%
            \textbf{\color{primaryColor}\faSlidersH\ Ajustes UI Accesible (CeRA):} \hspace{2mm}
            \fAlBadgesAjustes%
        }
    \end{tcolorbox}
    \newpage
    \endgroup
}

% COMANDO PARA PLANTILLA EN BLANCO DE ALUMNADO
\newcommand{\fichaAlumnoBlanco}{%
    \fichaAlumno[%
        bannerTitulo={\faUserEdit\ FICHA DE PERSONA: ALUMNADO \textbar\ PLANTILLA DE TRABAJO},
        cita={[Cita textual que exprese cómo vive su rutina, qué le ayuda a anticipar y qué le genera ansiedad o rechazo].}%
    ]%
}

\ifdefined\isSubFile
    \expandafter\endinput
\fi

\begin{document}
\fichaAlumnoBlanco
\end{document}
 y luego invocar \fichaAlumno[...] o \fichaAlumnoBlanco
% ==============================================================================

\ifdefined\isSubFile
    % Si se incluye desde otro documento, los paquetes ya están cargados
\else
    \documentclass[10pt]{article}
    \usepackage[utf8]{inputenc}
    \usepackage[spanish,es-tabla]{babel}
    \usepackage[margin=0.8cm,top=0.6cm,bottom=0.6cm]{geometry}
    \usepackage{helvet}
    \renewcommand{\familydefault}{\sfdefault}
    \usepackage{microtype}
    \usepackage{graphicx}
    \usepackage{xcolor}
    \usepackage{tikz}
    \usepackage{enumitem}
    \usepackage{fontawesome5}
    \usepackage{tcolorbox}
    \tcbuselibrary{skins}
    \pagestyle{empty}
\fi

% Colores del sistema
\providecolor{primaryColor}{RGB}{20, 65, 105}     % Azul corporativo
\providecolor{secondaryColor}{RGB}{35, 115, 170}  % Azul medio accesible
\providecolor{tealColor}{RGB}{20, 130, 140}       % Verde azulado
\providecolor{greenColor}{RGB}{30, 135, 75}       % Verde bosque
\providecolor{accentColor}{RGB}{215, 105, 25}     % Naranja cálido
\providecolor{purpleColor}{RGB}{120, 60, 150}     % Púrpura suave
\providecolor{dangerColor}{RGB}{185, 45, 35}      % Rojo criticidad
\providecolor{bgCard}{RGB}{252, 253, 255}         % Fondo limpio
\providecolor{borderCard}{RGB}{215, 225, 235}     % Borde sutil
\providecolor{textDark}{RGB}{30, 40, 50}          % Texto principal oscuro
\providecolor{textMuted}{RGB}{95, 110, 125}       % Texto secundario

% Configuración de listas
\setlist[itemize]{leftmargin=3.0mm, itemsep=0.4pt, parsep=0pt, topsep=0.4pt, partopsep=0pt}

% Estilos de caja fija para Alumno
\tcbset{
    cardAlumnoBox/.style={
        enhanced,
        boxrule=0.65pt,
        arc=1.8mm,
        height=32.0mm,
        left=5.0pt, right=5.0pt, top=1.8pt, bottom=1.8pt,
        boxsep=0pt,
        colback=bgCard,
        colframe=borderCard,
        fonttitle=\bfseries\small,
        fontupper=\fontsize{8.5pt}{10.0pt}\selectfont,
    },
    cardAlumnoWideBox/.style={
        enhanced,
        boxrule=0.65pt,
        arc=1.8mm,
        height=22.5mm,
        left=6.0pt, right=6.0pt, top=2.0pt, bottom=2.0pt,
        boxsep=0pt,
        colback=bgCard,
        colframe=borderCard,
        fonttitle=\bfseries\small,
    }
}

% Macro para badges visuales (si no existe)
\providecommand{\badge}[2]{%
    \tikz[baseline=(b.base)] \node[fill=#1!12, text=#1!90!black, draw=#1!35, line width=0.4pt, rounded corners=1.2mm, inner sep=2.0pt, font=\bfseries\scriptsize] (b) {#2};%
}

% Definición de claves con pgfkeys para \fichaAlumno
\pgfkeys{
    /fichaAlumno/.is family, /fichaAlumno,
    default/.style = {
        bannerTitulo = {\faUserGraduate\ FICHA DE PERSONA: ALUMNADO (METODOLOGÍA DCU)},
        subtitulo = {Proyecto AUTOVIDA (26-175) \textbar\ Reto CeRA — Agenda Accesible PTVAL \textbar\ Asignatura MDA (UGR)},
        equipo = {\textbf{Equipo CeRA-MDA} (5 integrantes)\\\ Course 2026/2027 \textbar\ \textbf{Práctica 1}},
        nombre = {[Nombre y Apellidos del Alumno]},
        nombreCorto = {[ALUMNO]},
        edad = {[Edad]},
        icono = {\faUserGraduate},
        foto = {},
        badgeRol = {[Perfil NEAE / Aula]},
        badgeTarea = {Tareas Rotativas},
        centro = {CEE Purísima Concepción},
        cita = {[Cita textual representativa que exprese cómo vive su rutina, qué le ayuda a anticipar y qué le genera ansiedad o rechazo].},
        demographic = {
            \item \textbf{Edad y Género (Age / Gender):} [[Edad] años, [Género]].
            \item \textbf{Ubicación (Location):} [[Granada capital / Municipio]].
            \item \textbf{Perfil diagnóstico:} [[Diagnóstico cognitivo, motor o sensorial]].
        },
        computerExp = {
            \item \textbf{Habilidades informáticas:} [[Nivel táctil: autónomo, asistido]].
            \item \textbf{Años de experiencia:} [[Años de contacto con tablets]].
            \item \textbf{Hábitos de uso:} [[Apps frecuentes: vídeos, pictogramas, juegos]].
            \item \textbf{Autonomía digital:} [[Baja / Media / Asistida]].
        },
        productExp = {
            \item \textbf{Productos similares:} [[PECS en papel, velcro, apps previas]].
            \item \textbf{Limitaciones actuales:} [[Pérdida de material, rigidez, frustración]].
            \item \textbf{Tendencias de uso:} [[Receptividad a imágenes reales vs pictogramas]].
        },
        tasks = {
            \item \textbf{Tareas Primarias:} 1) [[Login accesible]] $\rightarrow$ 2) [[Ver tarea actual]] $\rightarrow$ 3) [[Ejecutar pasos guiados]] $\rightarrow$ 4) [[Confirmar finalización]].
            \item \textbf{Tareas Secundarias:} [[Ver siguiente tarea / pedir ayuda]].
        },
        domainKnowledge = {
            \item \textbf{Comprensión del área:} [[Orientación en el colegio, aulas y profesores]].
            \item \textbf{Comprensión temporal:} [[Visual (disco/reloj de arena), no numérica]].
            \item \textbf{Destreza procedimental:} [[Grado de autonomía en la tarea asignada]].
        },
        technology = {
            \item \textbf{Hardware:} [[Tablet Android / pulsador externo]].
            \item \textbf{Software:} [[Android / síntesis de voz]].
            \item \textbf{Herramientas habituales:} [[Conexión Wi-Fi / atril]].
        },
        attitudes = {
            \item \textbf{Preferencias:} [[Refuerzos sensoriales favoritos, limpieza visual]].
            \item \textbf{Miedo a la tecnología:} [[Reacción a errores, sonidos o bloqueos]].
            \item \textbf{Valores:} [[Sentimiento de logro, orgullo por autonomía]].
        },
        learningStyle = {
            \item \textbf{Perfil de aprendizaje:} [[\textbf{Visual} / \textbf{Auditivo} / \textbf{Kinestésico}]].
            \item \textbf{Canales óptimos:} [[Fotos reales / Pictogramas ARASAAC / Vídeos]].
            \item \textbf{Metodología:} [[Modelado guiado, repetición, ensayo y éxito]].
        },
        criticality = {
            \item \textbf{Consecuencias de un error:} \textit{Emocional:} [[¿Bloqueo, crisis de frustración, rechazo?]] \textit{Operativo:} [[¿Desorientación, interrupción de la tarea?]]
            \item \textbf{Tolerancia del software:} [[Acciones no destructivas, confirmación siempre asistida, botón SOS visible]].
        },
        escenario = {[[Narrar en 3-4 líneas una situación real: a qué hora usa la agenda, cómo anticipa su rutina con apoyos visuales y cómo supera el momento crítico logrando autonomía sin estrés.]]},
        badgesAjustes = {
            \badge{primaryColor}{\faKey\ Login: [[Método]]} \hspace{2mm}
            \badge{secondaryColor}{\faIcons\ [[ARASAAC / Fotos]]} \hspace{2mm}
            \badge{accentColor}{\faClock\ [[Temporizador]]} \hspace{2mm}
            \badge{purpleColor}{\faVolumeUp\ [[Voz TTS]]}%
        }
    },
    bannerTitulo/.store in = \fAlBannerTitulo,
    subtitulo/.store in = \fAlSubtitulo,
    equipo/.store in = \fAlEquipo,
    nombre/.store in = \fAlNombre,
    nombreCorto/.store in = \fAlNombreCorto,
    edad/.store in = \fAlEdad,
    icono/.store in = \fAlIcono,
    foto/.store in = \fAlFoto,
    badgeRol/.store in = \fAlBadgeRol,
    badgeTarea/.store in = \fAlBadgeTarea,
    centro/.store in = \fAlCentro,
    cita/.store in = \fAlCita,
    demographic/.store in = \fAlDemographic,
    computerExp/.store in = \fAlComputerExp,
    productExp/.store in = \fAlProductExp,
    tasks/.store in = \fAlTasks,
    domainKnowledge/.store in = \fAlDomainKnowledge,
    technology/.store in = \fAlTechnology,
    attitudes/.store in = \fAlAttitudes,
    learningStyle/.store in = \fAlLearningStyle,
    criticality/.store in = \fAlCriticality,
    escenario/.store in = \fAlEscenario,
    badgesAjustes/.store in = \fAlBadgesAjustes,
}

% COMANDO PRINCIPAL \fichaAlumno
\newcommand{\fichaAlumno}[1][]{%
    \begingroup
    \pgfkeys{/fichaAlumno, default, #1}%
    % 1. BANNER SUPERIOR
    \begin{tcolorbox}[
        enhanced,
        colback=primaryColor,
        colframe=primaryColor,
        arc=1.6mm,
        left=7pt, right=7pt, top=2.5pt, bottom=2.5pt,
        boxsep=0pt, coltext=white
    ]
        \begin{minipage}[c]{0.70\textwidth}
            {\normalsize\textbf{\fAlBannerTitulo}}\\[0.2mm]
            {\tiny \fAlSubtitulo}
        \end{minipage}\hfill
        \begin{minipage}[c]{0.29\textwidth}
            \raggedleft\tiny
            \fAlEquipo
        \end{minipage}
    \end{tcolorbox}
    \vspace{0.8mm}
    % 2. TARJETA DE IDENTIDAD DEL ALUMNO
    \begin{tcolorbox}[
        enhanced,
        colback=white,
        colframe=primaryColor!65,
        arc=2.0mm,
        boxrule=0.8pt,
        left=6pt, right=6pt, top=3.0pt, bottom=3.0pt,
        boxsep=0pt,
        drop shadow=black!5
    ]
        \begin{minipage}[c]{0.10\textwidth}
            \centering
            \begin{tikzpicture}
                \ifx\fAlFoto\empty
                    \draw[fill=primaryColor!8, draw=primaryColor!30, line width=0.5pt, rounded corners=2mm] (0,0) rectangle (1.4, 1.55);
                    \node at (0.7, 0.95) {\color{primaryColor}\scalebox{1.5}{\fAlIcono}};
                    \node[align=center] at (0.7, 0.28) {\tiny\textbf{\color{primaryColor}\fAlNombreCorto}\\\fontsize{4pt}{5pt}\selectfont\color{textMuted}\fAlEdad};
                \else
                    \begin{scope}
                        \clip[rounded corners=2mm] (0,0) rectangle (1.45, 1.55);
                        \node[inner sep=0pt] at (0.725, 0.775) {\includegraphics[width=1.45cm,height=1.55cm,keepaspectratio=false]{\fAlFoto}};
                    \end{scope}
                    \draw[draw=primaryColor!50, line width=0.8pt, rounded corners=2mm] (0,0) rectangle (1.45, 1.55);
                \fi
            \end{tikzpicture}
        \end{minipage}\hfill
        \begin{minipage}[c]{0.88\textwidth}
            \begin{minipage}[b]{0.68\textwidth}\raggedright
                {\large\textbf{\color{primaryColor}\fAlNombre}} \hspace{2mm}
                \badge{secondaryColor}{\fAlBadgeRol} \badge{accentColor}{\fAlBadgeTarea}
            \end{minipage}\hfill
            \begin{minipage}[b]{0.31\textwidth}
                \raggedleft\small\textbf{\color{primaryColor}\faSchool\ \fAlCentro}
            \end{minipage}\\[0.6mm]
            \tikz[baseline] \node[fill=primaryColor!5, draw=primaryColor!25, rounded corners=1.2mm, inner sep=2.5pt, text width=0.98\linewidth] {
                \small\itshape\color{primaryColor}
                \faQuoteLeft\ ``\fAlCita''\ \faQuoteRight
            };
        \end{minipage}
    \end{tcolorbox}
    \vspace{0.8mm}
    % 3. LAS 9 CAJAS DE ALUMNOS (4 FILAS X 2 COLS + 1 FILA COMPLETA)
    \noindent\begin{minipage}[t]{0.485\textwidth}
        \begin{tcolorbox}[cardAlumnoBox, colframe=primaryColor!50, coltitle=primaryColor, title={\faIdCard\ 1. Demographic characteristics}]
            \fontsize{8.2pt}{9.8pt}\selectfont
            \begin{itemize}
                \fAlDemographic
            \end{itemize}
        \end{tcolorbox}
    \end{minipage}\hfill
    \begin{minipage}[t]{0.485\textwidth}
        \begin{tcolorbox}[cardAlumnoBox, colframe=secondaryColor!60, coltitle=secondaryColor!90!black, title={\faLaptop\ 2. Computer experience}]
            \fontsize{8.2pt}{9.8pt}\selectfont
            \begin{itemize}
                \fAlComputerExp
            \end{itemize}
        \end{tcolorbox}
    \end{minipage}
    \vspace{0.8mm}
    \noindent\begin{minipage}[t]{0.485\textwidth}
        \begin{tcolorbox}[cardAlumnoBox, colframe=secondaryColor!60, coltitle=secondaryColor!90!black, title={\faSearchPlus\ 3. Specific product experience}]
            \fontsize{8.2pt}{9.8pt}\selectfont
            \begin{itemize}
                \fAlProductExp
            \end{itemize}
        \end{tcolorbox}
    \end{minipage}\hfill
    \begin{minipage}[t]{0.485\textwidth}
        \begin{tcolorbox}[cardAlumnoBox, colframe=greenColor!65, coltitle=greenColor!90!black, title={\faTasks\ 4. Tasks (Primary and secondary)}]
            \fontsize{8.2pt}{9.8pt}\selectfont
            \begin{itemize}
                \fAlTasks
            \end{itemize}
        \end{tcolorbox}
    \end{minipage}
    \vspace{0.8mm}
    \noindent\begin{minipage}[t]{0.485\textwidth}
        \begin{tcolorbox}[cardAlumnoBox, colframe=accentColor!65, coltitle=accentColor!90!black, title={\faBrain\ 5. Domain knowledge}]
            \fontsize{8.2pt}{9.8pt}\selectfont
            \begin{itemize}
                \fAlDomainKnowledge
            \end{itemize}
        \end{tcolorbox}
    \end{minipage}\hfill
    \begin{minipage}[t]{0.485\textwidth}
        \begin{tcolorbox}[cardAlumnoBox, colframe=tealColor!60, coltitle=tealColor!90!black, title={\faTools\ 6. Technology available}]
            \fontsize{8.2pt}{9.8pt}\selectfont
            \begin{itemize}
                \fAlTechnology
            \end{itemize}
        \end{tcolorbox}
    \end{minipage}
    \vspace{0.8mm}
    \noindent\begin{minipage}[t]{0.485\textwidth}
        \begin{tcolorbox}[cardAlumnoBox, colframe=purpleColor!60, coltitle=purpleColor!90!black, title={\faHeart\ 7. Attitudes and values}]
            \fontsize{8.2pt}{9.8pt}\selectfont
            \begin{itemize}
                \fAlAttitudes
            \end{itemize}
        \end{tcolorbox}
    \end{minipage}\hfill
    \begin{minipage}[t]{0.485\textwidth}
        \begin{tcolorbox}[cardAlumnoBox, colframe=purpleColor!60, coltitle=purpleColor!90!black, title={\faEye\ 8. Learning style}]
            \fontsize{8.2pt}{9.8pt}\selectfont
            \begin{itemize}
                \fAlLearningStyle
            \end{itemize}
        \end{tcolorbox}
    \end{minipage}
    \vspace{0.8mm}
    % FILA 5: Punto 9 (ANCHO COMPLETO — CRITICALITY OF ERRORS)
    \noindent\begin{minipage}[t]{\textwidth}
        \begin{tcolorbox}[cardAlumnoWideBox, colframe=dangerColor!65, coltitle=dangerColor!90!black, title={\faExclamationTriangle\ 9. Criticality of errors (Consecuencias de un error y tolerancia del software)}]
            \fontsize{8.2pt}{9.8pt}\selectfont
            \begin{itemize}
                \fAlCriticality
            \end{itemize}
        \end{tcolorbox}
    \end{minipage}
    \vspace{0.8mm}
    % 4. BARRA INFERIOR DE INTEGRACIÓN CeRA
    \begin{tcolorbox}[
        enhanced,
        colback=primaryColor!4,
        colframe=primaryColor!45,
        arc=1.6mm,
        left=6pt, right=6pt, top=2.5pt, bottom=2.5pt,
        boxsep=0pt
    ]
        \scriptsize
        \textbf{\color{primaryColor}\faMapSigns\ Escenario de Uso Contextual (Día típico):} \fAlEscenario\\
        \vspace{0.3mm}
        \hrule height 0.35pt \relax
        \vspace{0.5mm}
        \centerline{%
            \textbf{\color{primaryColor}\faSlidersH\ Ajustes UI Accesible (CeRA):} \hspace{2mm}
            \fAlBadgesAjustes%
        }
    \end{tcolorbox}
    \newpage
    \endgroup
}

% COMANDO PARA PLANTILLA EN BLANCO DE ALUMNADO
\newcommand{\fichaAlumnoBlanco}{%
    \fichaAlumno[%
        bannerTitulo={\faUserEdit\ FICHA DE PERSONA: ALUMNADO \textbar\ PLANTILLA DE TRABAJO},
        cita={[Cita textual que exprese cómo vive su rutina, qué le ayuda a anticipar y qué le genera ansiedad o rechazo].}%
    ]%
}

\ifdefined\isSubFile
    \expandafter\endinput
\fi

\begin{document}
\fichaAlumnoBlanco
\end{document}
 y luego invocar \fichaAlumno[...] o \fichaAlumnoBlanco
% ==============================================================================

\ifdefined\isSubFile
    % Si se incluye desde otro documento, los paquetes ya están cargados
\else
    \documentclass[10pt]{article}
    \usepackage[utf8]{inputenc}
    \usepackage[spanish,es-tabla]{babel}
    \usepackage[margin=0.8cm,top=0.6cm,bottom=0.6cm]{geometry}
    \usepackage{helvet}
    \renewcommand{\familydefault}{\sfdefault}
    \usepackage{microtype}
    \usepackage{graphicx}
    \usepackage{xcolor}
    \usepackage{tikz}
    \usepackage{enumitem}
    \usepackage{fontawesome5}
    \usepackage{tcolorbox}
    \tcbuselibrary{skins}
    \pagestyle{empty}
\fi

% Colores del sistema
\providecolor{primaryColor}{RGB}{20, 65, 105}     % Azul corporativo
\providecolor{secondaryColor}{RGB}{35, 115, 170}  % Azul medio accesible
\providecolor{tealColor}{RGB}{20, 130, 140}       % Verde azulado
\providecolor{greenColor}{RGB}{30, 135, 75}       % Verde bosque
\providecolor{accentColor}{RGB}{215, 105, 25}     % Naranja cálido
\providecolor{purpleColor}{RGB}{120, 60, 150}     % Púrpura suave
\providecolor{dangerColor}{RGB}{185, 45, 35}      % Rojo criticidad
\providecolor{bgCard}{RGB}{252, 253, 255}         % Fondo limpio
\providecolor{borderCard}{RGB}{215, 225, 235}     % Borde sutil
\providecolor{textDark}{RGB}{30, 40, 50}          % Texto principal oscuro
\providecolor{textMuted}{RGB}{95, 110, 125}       % Texto secundario

% Configuración de listas
\setlist[itemize]{leftmargin=3.0mm, itemsep=0.4pt, parsep=0pt, topsep=0.4pt, partopsep=0pt}

% Estilos de caja fija para Alumno
\tcbset{
    cardAlumnoBox/.style={
        enhanced,
        boxrule=0.65pt,
        arc=1.8mm,
        height=32.0mm,
        left=5.0pt, right=5.0pt, top=1.8pt, bottom=1.8pt,
        boxsep=0pt,
        colback=bgCard,
        colframe=borderCard,
        fonttitle=\bfseries\small,
        fontupper=\fontsize{8.5pt}{10.0pt}\selectfont,
    },
    cardAlumnoWideBox/.style={
        enhanced,
        boxrule=0.65pt,
        arc=1.8mm,
        height=22.5mm,
        left=6.0pt, right=6.0pt, top=2.0pt, bottom=2.0pt,
        boxsep=0pt,
        colback=bgCard,
        colframe=borderCard,
        fonttitle=\bfseries\small,
    }
}

% Macro para badges visuales (si no existe)
\providecommand{\badge}[2]{%
    \tikz[baseline=(b.base)] \node[fill=#1!12, text=#1!90!black, draw=#1!35, line width=0.4pt, rounded corners=1.2mm, inner sep=2.0pt, font=\bfseries\scriptsize] (b) {#2};%
}

% Definición de claves con pgfkeys para \fichaAlumno
\pgfkeys{
    /fichaAlumno/.is family, /fichaAlumno,
    default/.style = {
        bannerTitulo = {\faUserGraduate\ FICHA DE PERSONA: ALUMNADO (METODOLOGÍA DCU)},
        subtitulo = {Proyecto AUTOVIDA (26-175) \textbar\ Reto CeRA — Agenda Accesible PTVAL \textbar\ Asignatura MDA (UGR)},
        equipo = {\textbf{Equipo CeRA-MDA} (5 integrantes)\\\ Course 2026/2027 \textbar\ \textbf{Práctica 1}},
        nombre = {[Nombre y Apellidos del Alumno]},
        nombreCorto = {[ALUMNO]},
        edad = {[Edad]},
        icono = {\faUserGraduate},
        foto = {},
        badgeRol = {[Perfil NEAE / Aula]},
        badgeTarea = {Tareas Rotativas},
        centro = {CEE Purísima Concepción},
        cita = {[Cita textual representativa que exprese cómo vive su rutina, qué le ayuda a anticipar y qué le genera ansiedad o rechazo].},
        demographic = {
            \item \textbf{Edad y Género (Age / Gender):} [[Edad] años, [Género]].
            \item \textbf{Ubicación (Location):} [[Granada capital / Municipio]].
            \item \textbf{Perfil diagnóstico:} [[Diagnóstico cognitivo, motor o sensorial]].
        },
        computerExp = {
            \item \textbf{Habilidades informáticas:} [[Nivel táctil: autónomo, asistido]].
            \item \textbf{Años de experiencia:} [[Años de contacto con tablets]].
            \item \textbf{Hábitos de uso:} [[Apps frecuentes: vídeos, pictogramas, juegos]].
            \item \textbf{Autonomía digital:} [[Baja / Media / Asistida]].
        },
        productExp = {
            \item \textbf{Productos similares:} [[PECS en papel, velcro, apps previas]].
            \item \textbf{Limitaciones actuales:} [[Pérdida de material, rigidez, frustración]].
            \item \textbf{Tendencias de uso:} [[Receptividad a imágenes reales vs pictogramas]].
        },
        tasks = {
            \item \textbf{Tareas Primarias:} 1) [[Login accesible]] $\rightarrow$ 2) [[Ver tarea actual]] $\rightarrow$ 3) [[Ejecutar pasos guiados]] $\rightarrow$ 4) [[Confirmar finalización]].
            \item \textbf{Tareas Secundarias:} [[Ver siguiente tarea / pedir ayuda]].
        },
        domainKnowledge = {
            \item \textbf{Comprensión del área:} [[Orientación en el colegio, aulas y profesores]].
            \item \textbf{Comprensión temporal:} [[Visual (disco/reloj de arena), no numérica]].
            \item \textbf{Destreza procedimental:} [[Grado de autonomía en la tarea asignada]].
        },
        technology = {
            \item \textbf{Hardware:} [[Tablet Android / pulsador externo]].
            \item \textbf{Software:} [[Android / síntesis de voz]].
            \item \textbf{Herramientas habituales:} [[Conexión Wi-Fi / atril]].
        },
        attitudes = {
            \item \textbf{Preferencias:} [[Refuerzos sensoriales favoritos, limpieza visual]].
            \item \textbf{Miedo a la tecnología:} [[Reacción a errores, sonidos o bloqueos]].
            \item \textbf{Valores:} [[Sentimiento de logro, orgullo por autonomía]].
        },
        learningStyle = {
            \item \textbf{Perfil de aprendizaje:} [[\textbf{Visual} / \textbf{Auditivo} / \textbf{Kinestésico}]].
            \item \textbf{Canales óptimos:} [[Fotos reales / Pictogramas ARASAAC / Vídeos]].
            \item \textbf{Metodología:} [[Modelado guiado, repetición, ensayo y éxito]].
        },
        criticality = {
            \item \textbf{Consecuencias de un error:} \textit{Emocional:} [[¿Bloqueo, crisis de frustración, rechazo?]] \textit{Operativo:} [[¿Desorientación, interrupción de la tarea?]]
            \item \textbf{Tolerancia del software:} [[Acciones no destructivas, confirmación siempre asistida, botón SOS visible]].
        },
        escenario = {[[Narrar en 3-4 líneas una situación real: a qué hora usa la agenda, cómo anticipa su rutina con apoyos visuales y cómo supera el momento crítico logrando autonomía sin estrés.]]},
        badgesAjustes = {
            \badge{primaryColor}{\faKey\ Login: [[Método]]} \hspace{2mm}
            \badge{secondaryColor}{\faIcons\ [[ARASAAC / Fotos]]} \hspace{2mm}
            \badge{accentColor}{\faClock\ [[Temporizador]]} \hspace{2mm}
            \badge{purpleColor}{\faVolumeUp\ [[Voz TTS]]}%
        }
    },
    bannerTitulo/.store in = \fAlBannerTitulo,
    subtitulo/.store in = \fAlSubtitulo,
    equipo/.store in = \fAlEquipo,
    nombre/.store in = \fAlNombre,
    nombreCorto/.store in = \fAlNombreCorto,
    edad/.store in = \fAlEdad,
    icono/.store in = \fAlIcono,
    foto/.store in = \fAlFoto,
    badgeRol/.store in = \fAlBadgeRol,
    badgeTarea/.store in = \fAlBadgeTarea,
    centro/.store in = \fAlCentro,
    cita/.store in = \fAlCita,
    demographic/.store in = \fAlDemographic,
    computerExp/.store in = \fAlComputerExp,
    productExp/.store in = \fAlProductExp,
    tasks/.store in = \fAlTasks,
    domainKnowledge/.store in = \fAlDomainKnowledge,
    technology/.store in = \fAlTechnology,
    attitudes/.store in = \fAlAttitudes,
    learningStyle/.store in = \fAlLearningStyle,
    criticality/.store in = \fAlCriticality,
    escenario/.store in = \fAlEscenario,
    badgesAjustes/.store in = \fAlBadgesAjustes,
}

% COMANDO PRINCIPAL \fichaAlumno
\newcommand{\fichaAlumno}[1][]{%
    \begingroup
    \pgfkeys{/fichaAlumno, default, #1}%
    % 1. BANNER SUPERIOR
    \begin{tcolorbox}[
        enhanced,
        colback=primaryColor,
        colframe=primaryColor,
        arc=1.6mm,
        left=7pt, right=7pt, top=2.5pt, bottom=2.5pt,
        boxsep=0pt, coltext=white
    ]
        \begin{minipage}[c]{0.70\textwidth}
            {\normalsize\textbf{\fAlBannerTitulo}}\\[0.2mm]
            {\tiny \fAlSubtitulo}
        \end{minipage}\hfill
        \begin{minipage}[c]{0.29\textwidth}
            \raggedleft\tiny
            \fAlEquipo
        \end{minipage}
    \end{tcolorbox}
    \vspace{0.8mm}
    % 2. TARJETA DE IDENTIDAD DEL ALUMNO
    \begin{tcolorbox}[
        enhanced,
        colback=white,
        colframe=primaryColor!65,
        arc=2.0mm,
        boxrule=0.8pt,
        left=6pt, right=6pt, top=3.0pt, bottom=3.0pt,
        boxsep=0pt,
        drop shadow=black!5
    ]
        \begin{minipage}[c]{0.10\textwidth}
            \centering
            \begin{tikzpicture}
                \ifx\fAlFoto\empty
                    \draw[fill=primaryColor!8, draw=primaryColor!30, line width=0.5pt, rounded corners=2mm] (0,0) rectangle (1.4, 1.55);
                    \node at (0.7, 0.95) {\color{primaryColor}\scalebox{1.5}{\fAlIcono}};
                    \node[align=center] at (0.7, 0.28) {\tiny\textbf{\color{primaryColor}\fAlNombreCorto}\\\fontsize{4pt}{5pt}\selectfont\color{textMuted}\fAlEdad};
                \else
                    \begin{scope}
                        \clip[rounded corners=2mm] (0,0) rectangle (1.45, 1.55);
                        \node[inner sep=0pt] at (0.725, 0.775) {\includegraphics[width=1.45cm,height=1.55cm,keepaspectratio=false]{\fAlFoto}};
                    \end{scope}
                    \draw[draw=primaryColor!50, line width=0.8pt, rounded corners=2mm] (0,0) rectangle (1.45, 1.55);
                \fi
            \end{tikzpicture}
        \end{minipage}\hfill
        \begin{minipage}[c]{0.88\textwidth}
            \begin{minipage}[b]{0.68\textwidth}\raggedright
                {\large\textbf{\color{primaryColor}\fAlNombre}} \hspace{2mm}
                \badge{secondaryColor}{\fAlBadgeRol} \badge{accentColor}{\fAlBadgeTarea}
            \end{minipage}\hfill
            \begin{minipage}[b]{0.31\textwidth}
                \raggedleft\small\textbf{\color{primaryColor}\faSchool\ \fAlCentro}
            \end{minipage}\\[0.6mm]
            \tikz[baseline] \node[fill=primaryColor!5, draw=primaryColor!25, rounded corners=1.2mm, inner sep=2.5pt, text width=0.98\linewidth] {
                \small\itshape\color{primaryColor}
                \faQuoteLeft\ ``\fAlCita''\ \faQuoteRight
            };
        \end{minipage}
    \end{tcolorbox}
    \vspace{0.8mm}
    % 3. LAS 9 CAJAS DE ALUMNOS (4 FILAS X 2 COLS + 1 FILA COMPLETA)
    \noindent\begin{minipage}[t]{0.485\textwidth}
        \begin{tcolorbox}[cardAlumnoBox, colframe=primaryColor!50, coltitle=primaryColor, title={\faIdCard\ 1. Demographic characteristics}]
            \fontsize{8.2pt}{9.8pt}\selectfont
            \begin{itemize}
                \fAlDemographic
            \end{itemize}
        \end{tcolorbox}
    \end{minipage}\hfill
    \begin{minipage}[t]{0.485\textwidth}
        \begin{tcolorbox}[cardAlumnoBox, colframe=secondaryColor!60, coltitle=secondaryColor!90!black, title={\faLaptop\ 2. Computer experience}]
            \fontsize{8.2pt}{9.8pt}\selectfont
            \begin{itemize}
                \fAlComputerExp
            \end{itemize}
        \end{tcolorbox}
    \end{minipage}
    \vspace{0.8mm}
    \noindent\begin{minipage}[t]{0.485\textwidth}
        \begin{tcolorbox}[cardAlumnoBox, colframe=secondaryColor!60, coltitle=secondaryColor!90!black, title={\faSearchPlus\ 3. Specific product experience}]
            \fontsize{8.2pt}{9.8pt}\selectfont
            \begin{itemize}
                \fAlProductExp
            \end{itemize}
        \end{tcolorbox}
    \end{minipage}\hfill
    \begin{minipage}[t]{0.485\textwidth}
        \begin{tcolorbox}[cardAlumnoBox, colframe=greenColor!65, coltitle=greenColor!90!black, title={\faTasks\ 4. Tasks (Primary and secondary)}]
            \fontsize{8.2pt}{9.8pt}\selectfont
            \begin{itemize}
                \fAlTasks
            \end{itemize}
        \end{tcolorbox}
    \end{minipage}
    \vspace{0.8mm}
    \noindent\begin{minipage}[t]{0.485\textwidth}
        \begin{tcolorbox}[cardAlumnoBox, colframe=accentColor!65, coltitle=accentColor!90!black, title={\faBrain\ 5. Domain knowledge}]
            \fontsize{8.2pt}{9.8pt}\selectfont
            \begin{itemize}
                \fAlDomainKnowledge
            \end{itemize}
        \end{tcolorbox}
    \end{minipage}\hfill
    \begin{minipage}[t]{0.485\textwidth}
        \begin{tcolorbox}[cardAlumnoBox, colframe=tealColor!60, coltitle=tealColor!90!black, title={\faTools\ 6. Technology available}]
            \fontsize{8.2pt}{9.8pt}\selectfont
            \begin{itemize}
                \fAlTechnology
            \end{itemize}
        \end{tcolorbox}
    \end{minipage}
    \vspace{0.8mm}
    \noindent\begin{minipage}[t]{0.485\textwidth}
        \begin{tcolorbox}[cardAlumnoBox, colframe=purpleColor!60, coltitle=purpleColor!90!black, title={\faHeart\ 7. Attitudes and values}]
            \fontsize{8.2pt}{9.8pt}\selectfont
            \begin{itemize}
                \fAlAttitudes
            \end{itemize}
        \end{tcolorbox}
    \end{minipage}\hfill
    \begin{minipage}[t]{0.485\textwidth}
        \begin{tcolorbox}[cardAlumnoBox, colframe=purpleColor!60, coltitle=purpleColor!90!black, title={\faEye\ 8. Learning style}]
            \fontsize{8.2pt}{9.8pt}\selectfont
            \begin{itemize}
                \fAlLearningStyle
            \end{itemize}
        \end{tcolorbox}
    \end{minipage}
    \vspace{0.8mm}
    % FILA 5: Punto 9 (ANCHO COMPLETO — CRITICALITY OF ERRORS)
    \noindent\begin{minipage}[t]{\textwidth}
        \begin{tcolorbox}[cardAlumnoWideBox, colframe=dangerColor!65, coltitle=dangerColor!90!black, title={\faExclamationTriangle\ 9. Criticality of errors (Consecuencias de un error y tolerancia del software)}]
            \fontsize{8.2pt}{9.8pt}\selectfont
            \begin{itemize}
                \fAlCriticality
            \end{itemize}
        \end{tcolorbox}
    \end{minipage}
    \vspace{0.8mm}
    % 4. BARRA INFERIOR DE INTEGRACIÓN CeRA
    \begin{tcolorbox}[
        enhanced,
        colback=primaryColor!4,
        colframe=primaryColor!45,
        arc=1.6mm,
        left=6pt, right=6pt, top=2.5pt, bottom=2.5pt,
        boxsep=0pt
    ]
        \scriptsize
        \textbf{\color{primaryColor}\faMapSigns\ Escenario de Uso Contextual (Día típico):} \fAlEscenario\\
        \vspace{0.3mm}
        \hrule height 0.35pt \relax
        \vspace{0.5mm}
        \centerline{%
            \textbf{\color{primaryColor}\faSlidersH\ Ajustes UI Accesible (CeRA):} \hspace{2mm}
            \fAlBadgesAjustes%
        }
    \end{tcolorbox}
    \newpage
    \endgroup
}

% COMANDO PARA PLANTILLA EN BLANCO DE ALUMNADO
\newcommand{\fichaAlumnoBlanco}{%
    \fichaAlumno[%
        bannerTitulo={\faUserEdit\ FICHA DE PERSONA: ALUMNADO \textbar\ PLANTILLA DE TRABAJO},
        cita={[Cita textual que exprese cómo vive su rutina, qué le ayuda a anticipar y qué le genera ansiedad o rechazo].}%
    ]%
}

\ifdefined\isSubFile
    \expandafter\endinput
\fi

\begin{document}
\fichaAlumnoBlanco
\end{document}
